%% kapa: exam=1 round=30 year=2019 period=2 form=A subject=accounting subject_ko=회계학 seq=8 number=48
%% source: Q-Net 감정평가사 기출문제 (cst00309) artlSeq=5205205 file=extracted/1008/5205205/2019년 제30회 감정평가사 1차시험 문제지/2019년 제30회 감정평가사 1차시험 문제지/2019년도 제30회 감정평가사 1차 2교시 A형.hwp
%% conversion: kordoc 4.24.51 -> tools/gen1.py (auto); review=proofread-sonnet; flags=none
\begin{question}{48}
(주)감평은 20x1년 초 기계장치(취득원가 ￦50,000, 내용연수 4년, 잔존가치 ￦0)를 취득하여 연수합계법으로 감가상각하고 있다. (주)감평은 20x1년 말 동 자산에 손상징후가 존재하여 회수가능액을 추정하였다. 그 결과 기계장치의 처분공정가치는 ￦25,000, 처분부대원가는 ￦3,000, 그리고 사용가치는 ￦23,000으로 확인되었다. (주)감평이 원가모형을 채택할 때, 동 기계장치와 관련하여 20x1년도에 인식할 손상차손은?
\begin{choices}
  \item ￦4,000
  \item ￦5,000
  \item ￦6,000
  \item ￦7,000
  \item ￦8,000
\end{choices}
\end{question}
\answer{4}
